% Options for packages loaded elsewhere
\PassOptionsToPackage{unicode}{hyperref}
\PassOptionsToPackage{hyphens}{url}
\PassOptionsToPackage{dvipsnames,svgnames,x11names}{xcolor}
%
\documentclass[
  10pt,
]{article}
\usepackage{amsmath,amssymb}
\usepackage{setspace}
\usepackage{iftex}
\ifPDFTeX
  \usepackage[T1]{fontenc}
  \usepackage[utf8]{inputenc}
  \usepackage{textcomp} % provide euro and other symbols
\else % if luatex or xetex
  \usepackage{unicode-math} % this also loads fontspec
  \defaultfontfeatures{Scale=MatchLowercase}
  \defaultfontfeatures[\rmfamily]{Ligatures=TeX,Scale=1}
\fi
\usepackage{lmodern}
\ifPDFTeX\else
  % xetex/luatex font selection
\fi
% Use upquote if available, for straight quotes in verbatim environments
\IfFileExists{upquote.sty}{\usepackage{upquote}}{}
\IfFileExists{microtype.sty}{% use microtype if available
  \usepackage[]{microtype}
  \UseMicrotypeSet[protrusion]{basicmath} % disable protrusion for tt fonts
}{}
\makeatletter
\@ifundefined{KOMAClassName}{% if non-KOMA class
  \IfFileExists{parskip.sty}{%
    \usepackage{parskip}
  }{% else
    \setlength{\parindent}{0pt}
    \setlength{\parskip}{6pt plus 2pt minus 1pt}}
}{% if KOMA class
  \KOMAoptions{parskip=half}}
\makeatother
\usepackage{xcolor}
\usepackage[a4paper,margin=27mm]{geometry}
\usepackage{color}
\usepackage{fancyvrb}
\newcommand{\VerbBar}{|}
\newcommand{\VERB}{\Verb[commandchars=\\\{\}]}
\DefineVerbatimEnvironment{Highlighting}{Verbatim}{commandchars=\\\{\}}
% Add ',fontsize=\small' for more characters per line
\newenvironment{Shaded}{}{}
\newcommand{\AlertTok}[1]{\textcolor[rgb]{1.00,0.00,0.00}{\textbf{#1}}}
\newcommand{\AnnotationTok}[1]{\textcolor[rgb]{0.38,0.63,0.69}{\textbf{\textit{#1}}}}
\newcommand{\AttributeTok}[1]{\textcolor[rgb]{0.49,0.56,0.16}{#1}}
\newcommand{\BaseNTok}[1]{\textcolor[rgb]{0.25,0.63,0.44}{#1}}
\newcommand{\BuiltInTok}[1]{\textcolor[rgb]{0.00,0.50,0.00}{#1}}
\newcommand{\CharTok}[1]{\textcolor[rgb]{0.25,0.44,0.63}{#1}}
\newcommand{\CommentTok}[1]{\textcolor[rgb]{0.38,0.63,0.69}{\textit{#1}}}
\newcommand{\CommentVarTok}[1]{\textcolor[rgb]{0.38,0.63,0.69}{\textbf{\textit{#1}}}}
\newcommand{\ConstantTok}[1]{\textcolor[rgb]{0.53,0.00,0.00}{#1}}
\newcommand{\ControlFlowTok}[1]{\textcolor[rgb]{0.00,0.44,0.13}{\textbf{#1}}}
\newcommand{\DataTypeTok}[1]{\textcolor[rgb]{0.56,0.13,0.00}{#1}}
\newcommand{\DecValTok}[1]{\textcolor[rgb]{0.25,0.63,0.44}{#1}}
\newcommand{\DocumentationTok}[1]{\textcolor[rgb]{0.73,0.13,0.13}{\textit{#1}}}
\newcommand{\ErrorTok}[1]{\textcolor[rgb]{1.00,0.00,0.00}{\textbf{#1}}}
\newcommand{\ExtensionTok}[1]{#1}
\newcommand{\FloatTok}[1]{\textcolor[rgb]{0.25,0.63,0.44}{#1}}
\newcommand{\FunctionTok}[1]{\textcolor[rgb]{0.02,0.16,0.49}{#1}}
\newcommand{\ImportTok}[1]{\textcolor[rgb]{0.00,0.50,0.00}{\textbf{#1}}}
\newcommand{\InformationTok}[1]{\textcolor[rgb]{0.38,0.63,0.69}{\textbf{\textit{#1}}}}
\newcommand{\KeywordTok}[1]{\textcolor[rgb]{0.00,0.44,0.13}{\textbf{#1}}}
\newcommand{\NormalTok}[1]{#1}
\newcommand{\OperatorTok}[1]{\textcolor[rgb]{0.40,0.40,0.40}{#1}}
\newcommand{\OtherTok}[1]{\textcolor[rgb]{0.00,0.44,0.13}{#1}}
\newcommand{\PreprocessorTok}[1]{\textcolor[rgb]{0.74,0.48,0.00}{#1}}
\newcommand{\RegionMarkerTok}[1]{#1}
\newcommand{\SpecialCharTok}[1]{\textcolor[rgb]{0.25,0.44,0.63}{#1}}
\newcommand{\SpecialStringTok}[1]{\textcolor[rgb]{0.73,0.40,0.53}{#1}}
\newcommand{\StringTok}[1]{\textcolor[rgb]{0.25,0.44,0.63}{#1}}
\newcommand{\VariableTok}[1]{\textcolor[rgb]{0.10,0.09,0.49}{#1}}
\newcommand{\VerbatimStringTok}[1]{\textcolor[rgb]{0.25,0.44,0.63}{#1}}
\newcommand{\WarningTok}[1]{\textcolor[rgb]{0.38,0.63,0.69}{\textbf{\textit{#1}}}}
\usepackage{longtable,booktabs,array}
\usepackage{calc} % for calculating minipage widths
% Correct order of tables after \paragraph or \subparagraph
\usepackage{etoolbox}
\makeatletter
\patchcmd\longtable{\par}{\if@noskipsec\mbox{}\fi\par}{}{}
\makeatother
% Allow footnotes in longtable head/foot
\IfFileExists{footnotehyper.sty}{\usepackage{footnotehyper}}{\usepackage{footnote}}
\makesavenoteenv{longtable}
\setlength{\emergencystretch}{3em} % prevent overfull lines
\providecommand{\tightlist}{%
  \setlength{\itemsep}{0pt}\setlength{\parskip}{0pt}}
\setcounter{secnumdepth}{-\maxdimen} % remove section numbering
\usepackage{microtype}
\usepackage{fancyhdr}
\usepackage{enumitem}
\usepackage{xcolor}
\usepackage{titlesec}
\definecolor{DCPBlue}{HTML}{204A74}
\definecolor{DCPGray}{HTML}{4A4A4A}
\pagestyle{fancy}
\fancyhf{}
\lhead{\small\color{DCPGray}Author Technical Correction}
\rhead{\small\color{DCPGray}DCP Challenge}
\cfoot{\small\thepage}
\renewcommand{\headrulewidth}{0.3pt}
\setlength{\headheight}{14pt}
\setlength{\parindent}{0pt}
\setlength{\parskip}{0.55em plus 0.12em minus 0.08em}
\setlist{nosep,leftmargin=1.5em}
\titleformat{\section}{\Large\bfseries\color{DCPBlue}}{}{0pt}{}
\titleformat{\subsection}{\large\bfseries\color{DCPBlue}}{}{0pt}{}
\titleformat{\subsubsection}{\normalsize\bfseries\color{DCPBlue}}{}{0pt}{}
\ifLuaTeX
  \usepackage{selnolig}  % disable illegal ligatures
\fi
\usepackage{bookmark}
\IfFileExists{xurl.sty}{\usepackage{xurl}}{} % add URL line breaks if available
\urlstyle{same}
\hypersetup{
  pdftitle={Author Technical Correction and Updated Analysis},
  pdfauthor={Ruge Lin},
  colorlinks=true,
  linkcolor={DCPBlue},
  filecolor={Maroon},
  citecolor={Blue},
  urlcolor={DCPBlue},
  pdfcreator={LaTeX via pandoc}}

\title{Author Technical Correction and Updated Analysis}
\usepackage{etoolbox}
\makeatletter
\providecommand{\subtitle}[1]{% add subtitle to \maketitle
  \apptocmd{\@title}{\par {\large #1 \par}}{}{}
}
\makeatother
\subtitle{For `A quantum computation capability verification protocol
for NISQ devices with dihedral coset problem'}
\author{Ruge Lin}
\date{17 August 2026}

\begin{document}
\maketitle

{
\hypersetup{linkcolor=}
\setcounter{tocdepth}{2}
\tableofcontents
}
\setstretch{1.04}
\begin{quote}
\textbf{Status.} This is an author-maintained technical correction
hosted with the code repository. It is not an APS Erratum, has not been
peer reviewed as a replacement article, and does not alter the journal
version. The statements here are made by Ruge Lin and should not be read
as a statement on behalf of the journal or as an assertion of
endorsement by the coauthor.
\end{quote}

\subsection{Abstract}\label{abstract}

The 2022 article proposed a DCP-derived circuit, called ParitySolve, and
interpreted its success probability relative to a quantity denoted by
\(p_B\) as a verification of quantum-computation capability. A later
reanalysis shows that the circuit-level parity extraction is correct,
but the verification interpretation is not. The quantity \(p_B\) is the
success probability of one specified all-Hadamard decoder; it is neither
proved nor valid as an upper bound over all strategies using the same
product measurements and classical postprocessing. An exact two-sample
counterexample gives success \(25/32\), compared with \(23/32\) for the
published decoder, without adding any quantum operation. Consequently,
the acceptance condition \(p>p_B\) does not establish the claimed
capability.

This note records what survives, supplies an exact formula for the ideal
ParitySolve probability, corrects the numerical interpretation of the
parameter choices used in Figure 5, replaces the reported fluctuation
formula by the Bernoulli standard error, and gives the strongest
aggregate reconstruction supported by the IBM counts retained in the
repository. The defensible surviving interpretation is a structured
DCP-derived circuit workload, not an adversarially sound verification
protocol.

\subsection{1. Scope and bibliographic
record}\label{scope-and-bibliographic-record}

This note concerns:

Ruge Lin and Weiqiang Wen, ``A quantum computation capability
verification protocol for NISQ devices with dihedral coset problem,''
\emph{Physical Review A} \textbf{106}, 012430 (2022), DOI:
\texttt{10.1103/PhysRevA.106.012430}, arXiv:2202.06984.

The original repository state is preserved by the immutable GitHub
release and tag \texttt{paper-2022-original}. The present repository
branch adds new files but leaves the six original Python scripts
unchanged.

The correction has four purposes:

\begin{enumerate}
\def\labelenumi{\arabic{enumi}.}
\tightlist
\item
  preserve the valid circuit construction;
\item
  withdraw the unsound interpretation of \(p_B\);
\item
  make the corrected calculations exactly reproducible;
\item
  separate a later DCP-derived witness from the claims of the 2022
  article.
\end{enumerate}

Citation history, peer-review history, and the reception of the article
are outside the scope of this note.

\subsection{2. Circuit result that remains
correct}\label{circuit-result-that-remains-correct}

Let \(N=2^n\), and let a DCP sample be

\[
\lvert\psi_{x,s}\rangle
=
\frac{
\lvert0\rangle\lvert x\rangle
+
\lvert1\rangle\lvert x+s\bmod N\rangle
}{\sqrt2},
\qquad x,s\in\mathbb Z_N.
\]

After applying the Fourier transform \(F_N\) to the rotation register
and measuring the Fourier label \(k\), the reflection qubit is, up to a
global phase,

\[
\lvert\phi_{k,s}\rangle
=
\frac{\lvert0\rangle+\omega_N^{ks}\lvert1\rangle}{\sqrt2},
\qquad
\omega_N=e^{2\pi i/N}.
\]

Suppose two measured labels satisfy

\[
k_1-k_2\equiv \frac N2\pmod N.
\]

On the selected output branch of the two-reflection-qubit CNOT
operation, the remaining relative phase is

\[
\omega_N^{(k_1-k_2)s}
=
\omega_N^{Ns/2}
=
(-1)^s.
\]

A final Hadamard measurement therefore returns \(s\bmod2\) exactly in
the ideal model. This ParitySolve algebra is correct. The circuit
remains a meaningful workload involving DCP preparation, Fourier
transforms, measurement, collision search, inter-cell interference, and
readout.

The correction does not withdraw this circuit identity.

\subsection{3. The original baseline and its
limitation}\label{the-original-baseline-and-its-limitation}

\subsubsection{\texorpdfstring{3.1 What \(p_B\) actually
computes}{3.1 What p\_B actually computes}}\label{what-p_b-actually-computes}

The article considered applying Hadamard gates to every qubit of every
DCP sample. Write one measurement outcome as \((r,y)\), where
\(r\in\{0,1\}\) is the reflection result and \(y\in\mathbb Z_N\) is the
rotation-register result.

After averaging over the uniform preparation label \(x\), the exact
conditional probability is

\[
q_s(r,y)
=
\frac1{N^2}
\#\left\{
 x\in\mathbb Z_N:
 r=x\mathbin{\cdot}y
 \oplus
 ((x+s)\bmod N)\mathbin{\cdot}y
\right\},
\]

where the dot denotes the binary inner product.

For \(y=1\), only the least significant bit is selected, so

\[
r=x_0\oplus(x+s)_0=s\bmod2.
\]

Thus one sample yields a special parity-revealing outcome with
probability \(1/N\). The published decoder retained only this outcome
and guessed uniformly if none of the \(mt\) samples produced it. Its
success probability is

\[
\boxed{
 p_B
 =
 1-\frac12\left(\frac{N-1}{N}\right)^{mt}
}.
\]

This formula is correct for that specific decoder. It should therefore
be renamed descriptively, for example as the \emph{special-outcome all-H
decoder probability}.

\subsubsection{3.2 Why the one-sample calculation looked
convincing}\label{why-the-one-sample-calculation-looked-convincing}

For a single sample, the special-outcome decoder is optimal for
discriminating the parity after the all-H product measurement. The
problem is not the one-sample formula.

The problem is that all \(mt\) samples in one repetition share the same
hidden secret \(s\). The relevant parity-conditioned record is
consequently a mixture over fixed-secret product distributions,

\[
\frac2N
\sum_{s\equiv b\ (2)}
q_s^{\otimes mt},
\]

rather than the tensor power of a one-sample parity average. Outcomes
that are inconclusive individually can therefore be correlated through
the reused secret and become informative when processed together.

\subsubsection{3.3 Exact two-sample
counterexample}\label{exact-two-sample-counterexample}

Take

\[
N=4,\qquad m=2,\qquad t=1.
\]

Use exactly the same all-H product measurements as the published
baseline. The published decoder succeeds with

\[
p_B
=
1-\frac12\left(\frac34\right)^2
=
\frac{23}{32}.
\]

Now use the following classical decoder on the same two outcomes
\((r_1,y_1)\) and \((r_2,y_2)\):

\begin{enumerate}
\def\labelenumi{\arabic{enumi}.}
\tightlist
\item
  if either \(y_j=1\), return the corresponding \(r_j\);
\item
  otherwise, if \(y_1,y_2\in\{2,3\}\), return \(r_1\oplus r_2\);
\item
  otherwise, guess uniformly.
\end{enumerate}

The first event occurs with probability

\[
1-\left(\frac34\right)^2=\frac7{16}
\]

and reveals the parity exactly.

The additional event \(y_1,y_2\in\{2,3\}\) occurs with probability
\(1/4\), is disjoint from the first event, and gives conditional success
\(3/4\): for even \(s\), the XOR is always zero; for odd \(s\), it is
uniform. With equal parity priors, the decoder is therefore correct with
probability \(3/4\) on this event.

Hence

\[
\begin{aligned}
p_{\rm better}
&=
\frac7{16}
+
\frac14\cdot\frac34
+
\left(1-\frac7{16}-\frac14\right)\frac12\\[2mm]
&=
\frac{25}{32}.
\end{aligned}
\]

The improvement is

\[
\frac{25}{32}-\frac{23}{32}=\frac1{16}.
\]

No entangling operation, quantum memory, alternative measurement basis,
additional sample, or quantum postprocessing is used. The difference
comes solely from retaining correlations discarded by the original
classical decoder.

This counterexample disproves the interpretation of \(p_B\) as a general
upper bound for the intended weaker class. Therefore the implication

\[
p>p_B
\quad\Longrightarrow\quad
\text{verified quantum-computation capability}
\]

is withdrawn.

\subsubsection{3.4 Full likelihood
decoder}\label{full-likelihood-decoder}

For any fixed all-H record \(o_1,\ldots,o_\ell\), with \(\ell=mt\), the
Bayes-optimal parity likelihoods are

\[
L_b(o_1,\ldots,o_\ell)
=
\frac2N
\sum_{\substack{s=0\\s\equiv b\ (2)}}^{N-1}
\prod_{j=1}^{\ell}q_s(o_j),
\qquad b\in\{0,1\}.
\]

Returning the larger likelihood systematically uses all information in
the original measurement record. The repository implements this decoder
exactly for small cases and by fixed-seed Monte Carlo for the larger
published instances.

\subsection{4. Exact ideal ParitySolve
probability}\label{exact-ideal-paritysolve-probability}

The original analysis used upper and lower bounds for the probability of
finding a complementary Fourier-label collision. An exact finite formula
is available.

Partition the \(N\) labels into \(N/2\) complementary pairs

\[
\{a,a+N/2\}.
\]

Let \(k_{\rm nc}(N,m)\) be the probability that \(m\) independent
uniform labels contain no complementary pair. If exactly \(j\)
complementary pairs are occupied without a collision, then:

\begin{enumerate}
\def\labelenumi{\arabic{enumi}.}
\tightlist
\item
  choose the occupied pairs: \(\binom{N/2}{j}\);
\item
  choose one member of each pair: \(2^j\);
\item
  assign the \(m\) labelled samples surjectively to the \(j\) selected
  labels: \(j!\,{m\brace j}\).
\end{enumerate}

Therefore

\[
\boxed{
k_{\rm nc}(N,m)
=
\frac1{N^m}
\sum_{j=1}^{\min(m,N/2)}
\binom{N/2}{j}
2^j j!\,{m\brace j}
},
\]

where \({m\brace j}\) is a Stirling number of the second kind. The case
\(m=0\) is defined as \(k_{\rm nc}=1\).

In one batch, a complementary collision exists with probability
\(1-k_{\rm nc}\). The selected CNOT branch succeeds with probability
\(1/2\). Under the retry rule implemented in the original code, one
batch fails to return a definitive parity with probability

\[
\frac{1+k_{\rm nc}}2.
\]

After \(t\) batches, the exact ideal success probability, including a
final random guess, is

\[
\boxed{
p_{\rm exact}
=
1-\frac12
\left(\frac{1+k_{\rm nc}}2\right)^t
}.
\]

The repository checks the formula against exhaustive enumeration for all
test instances with \(n\le3\) and \(m\le4\).

\subsection{5. Corrected interpretation of the Figure 5 parameter
choices}\label{corrected-interpretation-of-the-figure-5-parameter-choices}

The original parameter search selected instances from the difference
between the published upper bound \(p_{\rm upper}\) and \(p_B\). That is
a legitimate bound comparison, but it does not establish the actual
honest separation. The corrected values are:

\begin{longtable}[]{@{}
  >{\raggedright\arraybackslash}p{(\columnwidth - 10\tabcolsep) * \real{0.1304}}
  >{\raggedleft\arraybackslash}p{(\columnwidth - 10\tabcolsep) * \real{0.1739}}
  >{\raggedleft\arraybackslash}p{(\columnwidth - 10\tabcolsep) * \real{0.1739}}
  >{\raggedleft\arraybackslash}p{(\columnwidth - 10\tabcolsep) * \real{0.1739}}
  >{\raggedleft\arraybackslash}p{(\columnwidth - 10\tabcolsep) * \real{0.1739}}
  >{\raggedleft\arraybackslash}p{(\columnwidth - 10\tabcolsep) * \real{0.1739}}@{}}
\toprule\noalign{}
\begin{minipage}[b]{\linewidth}\raggedright
Figure
\end{minipage} & \begin{minipage}[b]{\linewidth}\raggedleft
\((n,m,t)\)
\end{minipage} & \begin{minipage}[b]{\linewidth}\raggedleft
\(p_B\)
\end{minipage} & \begin{minipage}[b]{\linewidth}\raggedleft
Published \(p_{\rm upper}\)
\end{minipage} & \begin{minipage}[b]{\linewidth}\raggedleft
Exact honest \(p\)
\end{minipage} & \begin{minipage}[b]{\linewidth}\raggedleft
Full-likelihood all-H
\end{minipage} \\
\midrule\noalign{}
\endhead
\bottomrule\noalign{}
\endlastfoot
5(a) & \((4,6,1)\) & 0.660533 & 0.664083 & \textbf{0.652965} &
\textbf{0.810940}, exact \\
5(b) & \((6,9,4)\) & 0.716371 & 0.817392 & \textbf{0.810918} &
\textbf{0.937289}, Monte Carlo \\
5(c) & \((9,21,9)\) & 0.654461 & 0.906606 & \textbf{0.904804} &
\textbf{0.957183}, Monte Carlo \\
\end{longtable}

For Figure 5(a), the actual honest difference is negative:

\[
p_{\rm exact}-p_B=-0.007568.
\]

For Figure 5(b), the actual difference is

\[
0.094547,
\]

which is below 10 percent. Figure 5(c) retains an honest difference of
approximately 25.03 percent against the special-outcome decoder, but the
optimized decoder using the same all-H measurement data exceeds the
honest ParitySolve score. None of these comparisons supplies a sound
verification gap.

The Figure 5(a) optimized value is exact:

\[
\frac{27863673495}{34359738368}
\approx0.8109396293.
\]

The Figure 5(b) and 5(c) values are fixed-seed Monte Carlo estimates
with Wilson intervals recorded in
\texttt{results/figure5\_corrected.csv}.

\subsection{6. Corrected statistical
uncertainty}\label{corrected-statistical-uncertainty}

Each complete repetition returns one correct-or-incorrect bit. Let \(Y\)
be its correctness indicator. Irrespective of the internal decomposition
into definitive branches and random guesses,

\[
Y\sim\operatorname{Bernoulli}(p).
\]

For \(r\) independent repetitions, the empirical accuracy therefore has
standard error

\[
\boxed{
\operatorname{SE}(\widehat p)
=
\sqrt{\frac{p(1-p)}r}
}.
\]

The expression used in the original paper and scripts,

\[
\sqrt{\frac{1-p}{2r}},
\]

accounts for the random final guess conditional on failure but omits the
fluctuation in whether a definitive branch occurs.

At the exact Figure 5(b) value with \(r=1000\), the two formulas give

\[
0.009723
\quad\text{and}\quad
0.012383,
\]

respectively.

\subsection{7. IBM aggregate
reconstruction}\label{ibm-aggregate-reconstruction}

The article describes a postselection condition based on measured output
bits. In the original \texttt{IBM.py}, the prepared values \texttt{x0}
and \texttt{x1} are also used as the collision test through

\begin{Shaded}
\begin{Highlighting}[]
\ControlFlowTok{if}\NormalTok{ x0 }\OperatorTok{!=}\NormalTok{ x1:}
\end{Highlighting}
\end{Shaded}

The preparation labels are not the measured Fourier labels. This legacy
reconstruction consequently uses only cases B, C, F, and G in the
definitive branch, although all eight preparation cases contain selected
hardware counts. It also substitutes an idealized selection probability
for the observed selection rate.

Using all eight archived aggregate cases and their observed
selected-count rates gives

\[
p_{\rm IBM,aggregate}=0.7439808125.
\]

For comparison,

\[
p_{\rm legacy\ code}=0.7545878022,
\qquad
p_{\rm ideal}=0.7890625.
\]

A 100,000-trial multinomial bootstrap of the three retained aggregate
categories gives the indicative 95 percent interval

\[
[0.739820,\,0.748090].
\]

This interval measures finite-count uncertainty only. The repository
does not contain the raw shot records or timestamps, so temporal drift
and shot-level correlations cannot be reconstructed.

The defensible interpretation is that the archived data show nontrivial
fidelity in the selected Bell-type branch. They do not constitute a
direct end-to-end execution of the challenge, and Figure 9 should not be
regenerated from the legacy reconstruction logic as though it did.

\subsection{8. Effect on the
conclusions}\label{effect-on-the-conclusions}

\subsubsection{Retained}\label{retained}

\begin{itemize}
\tightlist
\item
  the DCP sample construction;
\item
  the post-Fourier reflection phase state;
\item
  the complementary-label collision mechanism;
\item
  the selected ParitySolve branch and its exact parity output;
\item
  the original collision expressions as upper and lower bounds;
\item
  the use of the circuit as a structured, hardware-sensitive workload;
\item
  the original clean/noisy simulations as illustrations of the specified
  circuit and noise model.
\end{itemize}

\subsubsection{Corrected or withdrawn}\label{corrected-or-withdrawn}

\begin{itemize}
\tightlist
\item
  \(p_B\) is not a general soundness bound;
\item
  \(p>p_B\) does not verify the claimed quantum-computation capability;
\item
  the statement that no analytic ideal probability is available is
  replaced by the exact formula above;
\item
  Figure 5 parameter gaps must be interpreted using the exact honest
  probability, not only an upper bound;
\item
  the standard-error formula is replaced by the Bernoulli formula;
\item
  the IBM reconstruction is replaced by the aggregate-only calculation
  above.
\end{itemize}

\subsubsection{Surviving scientific
classification}\label{surviving-scientific-classification}

The original construction is best classified as a \textbf{DCP-derived
circuit benchmark or workload}. Calling it an adversarial
capability-verification protocol would require an explicit null class
and a proved optimal soundness bound against that entire class. The 2022
threshold does not provide such a bound.

\subsection{9. Separate follow-up
result}\label{separate-follow-up-result}

The repository also contains a later phase-twirled, heralded DCP
construction under \texttt{witness/}. It is not a repair already present
in the 2022 article. Its ideal theorem concerns a narrower capability: a
nonseparable joint measurement across two trusted DCP input cells,
against POVMs whose effects are separable across the cell partition.

That follow-up is explicitly labelled non-peer-reviewed and is separated
from this correction. It does not restore the original claim of general
quantum-computation verification.

\subsection{10. Reproducibility}\label{reproducibility}

The decisive calculations can be reproduced from a modern Python
environment without Qibo:

\begin{Shaded}
\begin{Highlighting}[]
\ExtensionTok{python} \AttributeTok{{-}m}\NormalTok{ pip install }\AttributeTok{{-}e}\NormalTok{ .}\PreprocessorTok{[}\SpecialStringTok{test}\PreprocessorTok{]}
\ExtensionTok{python}\NormalTok{ examples/reproduce\_counterexample.py}
\ExtensionTok{python}\NormalTok{ examples/reproduce\_figure5.py}
\ExtensionTok{python}\NormalTok{ examples/reproduce\_ibm\_reanalysis.py}
\ExtensionTok{pytest} \AttributeTok{{-}q}
\end{Highlighting}
\end{Shaded}

The committed machine-readable outputs are:

\begin{itemize}
\tightlist
\item
  \texttt{results/correction\_core\_results.json};
\item
  \texttt{results/figure5\_corrected.csv};
\item
  \texttt{results/ibm\_aggregate\_reanalysis.json};
\item
  \texttt{results/heralded\_witness\_results.json}.
\end{itemize}

The six original scripts remain unchanged. They require their historical
software stack and are not imported by the corrected tests.

\subsection{Appendix A. Exact outcome table for the smallest
counterexample}\label{appendix-a.-exact-outcome-table-for-the-smallest-counterexample}

For \(N=4\), the integer counts \(16q_s(r,y)\), with columns ordered by
\((r=0,y=0,1,2,3)\) followed by \((r=1,y=0,1,2,3)\), are

\[
\begin{array}{c|rrrrrrrr}
s&0,0&0,1&0,2&0,3&1,0&1,1&1,2&1,3\\\hline
0&4&4&4&4&0&0&0&0\\
1&4&0&2&2&0&4&2&2\\
2&4&4&0&0&0&0&4&4\\
3&4&0&2&2&0&4&2&2
\end{array}
\]

This table makes both parts of the improved decoder transparent. The
\(y=1\) columns reveal the parity exactly. Conditional on
\(y\in\{2,3\}\), even secrets force a fixed reflection result while odd
secrets produce an unbiased result, which creates the two-sample
correlation.

\subsection{Appendix B. Repository
status}\label{appendix-b.-repository-status}

\begin{itemize}
\tightlist
\item
  Historical code snapshot: immutable tag and release
  \texttt{paper-2022-original}.
\item
  Author correction branch: \texttt{author-correction-2026}.
\item
  Original six Python scripts: preserved byte-for-byte.
\item
  Current note: author-maintained, not an APS Erratum.
\item
  Follow-up witness: separate, ideal-model, non-peer-reviewed.
\end{itemize}

\end{document}
